\documentclass[11pt, a4paper]{article}

\usepackage[a4paper, top=2.5cm, bottom=2.5cm, left=2cm, right=2cm]{geometry}
\usepackage{amsmath}
\usepackage{amssymb}
\usepackage[colorlinks=true, linkcolor=blue, urlcolor=blue, citecolor=blue]{hyperref}

\title{Theory of Everything - Volume 47 \\ \Large Transcendent Architectures}
\author{Omega Protocol}
\date{\today}

\begin{document}

\maketitle

\section*{Layman's Summary}
Once a civilization conquers its galaxy, what's next? Volume 47 explores Type IV civilizations—beings so advanced they don't just live in the universe; they rewrite its source code. We explore the mathematical mechanics of changing the laws of physics and creating entirely new, custom-built universes.

\section{Type IV Civilizations (Axiom 50)}
A Type I civilization harnesses a planet; a Type III harnesses a galaxy. We establish that a Type IV civilization is one capable of directly manipulating the underlying mathematical axioms (the Type III\_1 algebra) of its own universe. They no longer manipulate matter; they manipulate the geometry of existence itself.

\section{Parameter Tuning (Theorem)}
Can you change the speed of light? We prove that fundamental constants are not arbitrary numbers handed down by a creator; they are localized properties of the vacuum's informational state. An incredibly dense, advanced Archive node can theoretically "tune" these parameters locally by altering the relative entropy of the vacuum.

\section{Artificial Universes (Theorem)}
How do you build a new universe in a lab? We prove that creating a stable, disconnected daughter universe requires the controlled compression and "topological pinching" of a highly ordered scalar field. An advanced civilization could theoretically pinch off a piece of spacetime, encode it with custom laws of physics, and let it inflate into a completely new reality.

\end{document}
